% !PRINTED 643
\weberpage{643}

%<<CONT
therefore, when the substitution $x_h = \psi(\eta_h)$ is made for the $x$, they pass over into functions of the variables $a$ which likewise cannot vanish identically. By the theorem of \S.\,143, 1. one can therefore set for the quantities $a_0, a_1 \ldots a_{n-1}$ such rational numbers that the functions $(\varepsilon, x)$ and the differences $f_\alpha - f_\beta$, which then pass over into $(\varepsilon, \xi)$ and into $k_\alpha - k_\beta$, become different from zero, and that the $\xi_h$ also remain different from one another, which was to be proved.

Having thus established this, we introduce the system of variables $x_0, x_1 \ldots x_{n-1}$ and a system of variables $y$ depending on it,

\weq{5}{y_0 = \chi(x_0), \quad y_1 = \chi(x_1), \ldots y_{n-1} = \chi(x_{n-1})}

in which $\chi$ denotes the function (3), and set

\weq{6}{\dfrac{(\varepsilon^{r_\nu}, y)}{(\varepsilon^{r_\nu}, x)} = \Theta_\nu.}

If any permutation is applied to the indices of the $x$, the indices of the $y$ undergo the same permutation. If, then, we apply to (6) the cyclic permutation $s$, $\Theta_\nu$ does not change (by \S.\,184, 4.).

If, then, we represent $\Theta_\nu$ in the normal form of \S.\,183, its coefficients are cyclic functions of $x_0, x_1 \ldots x_{n-1}$.

If one applies to $\Theta_\nu$ the permutation $t^{-1} = (h, g^{-1}h)$, then $\Theta_\nu$ passes over, by \S.\,184, 4., into $\Theta_{\nu+1}$, or the $\Theta_\nu$ undergo a cyclic permutation.

But the substitution $\sigma = (\varepsilon, \varepsilon^g)$ has the same effect as well, and therefore the system of functions $\Theta_0, \Theta_1 \ldots \Theta_{n-2}$ has the properties $\alpha)\ \beta)\ \gamma)\ \delta)$ (\S.\,184).

By the theorem proved at the end of \S.\,184, therefore, when we replace the $x$ by the $\xi$ and consequently the $y$ by the $\eta$, $\Theta_\nu$ passes over into a rational function of $k_\nu$

\weq{7}{Q_\nu = Q(k_\nu)}

and from (6) it follows that

\weq{8}{(\varepsilon^{r_\nu}, \eta) = Q_\nu (\varepsilon^{r_\nu}, \xi).}

If one substitutes herein for $(\varepsilon^{r_\nu}, \xi)$ the expression (5), \S.\,185, one obtains for $(\varepsilon^{r_\nu}, \eta)$ an expression of entirely the same form, only with the difference that in place of $K_\nu$ there has entered $Q_\nu K_\nu$. The functions $Q_\nu K_\nu$ have essentially the same properties as the function $K_\nu$, only that they also
%CONT>>
